\documentclass{beamer}

\usepackage[francais]{babel}
\usepackage{rtxslides}
\usepackage{tikz}
\usetikzlibrary{shapes}

\title{\rtx\ -- Technical Assessment 2}
\author{Louis Opter \\ \texttt{<louis@lse.epitech.eu>}}

\begin{document}

\begin{frame}
\titlepage
\end{frame}

\begin{frame}
\rmfamily{\itshape{«~\rtx\ est un \textcolor{rathaxesred}{langage dédié} qui
permet de \textcolor{rathaxesred}{décrire des pilotes} de périphériques. \rtx\
compile vers des modules noyaux écrits en C pour \textcolor{rathaxesred}{Linux},
\textcolor{rathaxesred}{Windows} et \textcolor{rathaxesred}{OpenBSD}.~»}}
\end{frame}

\begin{frame}{Développement de pilotes}
\begin{itemize}
\item code critique ;
\item double compétence ;
\item code spécifique ;
\item auto formation.
\end{itemize}
\end{frame}

\begin{frame}{Objectifs de Rathaxes}
\begin{itemize}
\item imposer plus de vérifications ;
\item séparer les compétences ;
\item unifier le développement.
\end{itemize}
\end{frame}

\begin{frame}{Cas d'utilisations}

\begin{center}
\only<1>{\LARGE\itshape{{Trois cas d'utilisations :}}}

\only<2>{\begin{tikzpicture}
\draw (-4, 1.5) node[above, text width=2cm, text centered, rotate=45]{\small{Développeur de pilotes}} ;
\draw (-4, -1) node[below]{\includegraphics[width=2cm]{logos/nvidia_white}} ;
\draw (-4, 0) node[below]{\includegraphics[width=2cm]{logos/intel}} ;
\draw (-4, 0.5) node[below]{\includegraphics[width=2cm]{logos/broadcom_white}} ;
\draw (-4, -3) node{\ldots} ;

\draw (0, 1.5) node[above, text width=2cm, text centered, rotate=45]{\textcolor{black}{\small{Mainteneur Rathaxes}}} ;

\draw (4, 1.5) node[above, text width=2cm, text centered, rotate=45]{\textcolor{black}{\small{Développeur systèmes}}} ;
\fill[black] (4, 0) circle(1.13) ; % hidden placeholder to make the right overflow
\end{tikzpicture}}

\only<3>{\begin{tikzpicture}
\draw (-4, 1.5) node[above, text width=2cm, text centered, rotate=45]{\small{Développeur de pilotes}} ;
\draw (-4, -1) node[below]{\includegraphics[width=2cm]{logos/nvidia_white}} ;
\draw (-4, 0) node[below]{\includegraphics[width=2cm]{logos/intel}} ;
\draw (-4, 0.5) node[below]{\includegraphics[width=2cm]{logos/broadcom_white}} ;
\draw (-4, -3) node{\ldots} ;

\draw (0, 1.5) node[above, text width=2cm, text centered, rotate=45]{\textcolor{black}{\small{Mainteneur Rathaxes}}} ;

\draw (4, 1.5) node[above, text width=2cm, text centered, rotate=45]{\small{Développeur systèmes}} ;
\draw (4, 0.5) node[below]{\includegraphics[width=2cm]{logos/tux}} ;
\draw (4, 0) node[above]{\includegraphics[width=2cm]{logos/windows_7_white}} ;
\draw (4, -3) node{\ldots} ;
\end{tikzpicture}}

\only<4>{\begin{tikzpicture}
\draw (-4, 1.5) node[above, text width=2cm, text centered, rotate=45]{\small{Développeur de pilotes}} ;
\draw (-4, -1) node[below]{\includegraphics[width=2cm]{logos/nvidia_white}} ;
\draw (-4, 0) node[below]{\includegraphics[width=2cm]{logos/intel}} ;
\draw (-4, 0.5) node[below]{\includegraphics[width=2cm]{logos/broadcom_white}} ;
\draw (-4, -3) node{\ldots} ;

\draw (0, 1.5) node[above, text width=2cm, text centered, rotate=45]{\small{Mainteneur Rathaxes}} ;
\draw (-0.5, 0.5) node[below]{\includegraphics[width=2cm]{logo_latex}} ;

\draw (4, 1.5) node[above, text width=2cm, text centered, rotate=45]{\small{Développeur systèmes}} ;
\draw (4, 0.5) node[below]{\includegraphics[width=2cm]{logos/tux}} ;
\draw (4, 0) node[above]{\includegraphics[width=2cm]{logos/windows_7_white}} ;
\draw (4, -3) node{\ldots} ;
\end{tikzpicture}}
\end{center}

% Note: I just discovered that beamer support columns inside frame, I'll give
% it a try next time…

\transdissolve<2>
\transdissolve<3>
\transdissolve<4>
\end{frame}

% Note: I wanted to have only one frame for Architecture + Composants but then
% this first sentence is not properlly vertically centered.
\begin{frame}{Architecture}
\begin{center}
\LARGE{\itshape{Séparée en plusieurs composants :}}
\end{center}
\end{frame}

\begin{frame}{Composants}

\tikzset{componentarrow/.style={<-, >=latex, color=rathaxesred, ultra thick}}
\tikzset{component/.style={draw, rectangle, rounded corners=3pt, ultra thick, color=rathaxesred}}
\tikzset{hidden/.style={draw, rectangle, rounded corners=3pt, ultra thick, color=black}}

\begin{overprint}

\onslide<1>
\begin{center}
\begin{tikzpicture}
\tikzset{component/.style={draw, rectangle, rounded corners=3pt, ultra thick, color=rathaxesred}}
\node[component] (cnorm) at (-3, 0)  {\textcolor{white}{\LARGE{CNorm}}} ;
\node[hidden] (rtxparse) at (2, 2.1)     {\Large{RTX\_Parse}} ;
\node[hidden] (rtxtree)  at (2, 0.7)     {\Large{RTX\_Tree}} ;
\node[hidden] (rtxlink)  at (2, -0.7)    {\Large{RTX\_Link}} ;
\node[hidden] (rtxtpl)   at (2, -2.1)    {\Large{RTX\_Tpl}} ;
\node[hidden] (tools)    at (-5, -2.1)   {\Large{Tools}} ;
\end{tikzpicture}
\end{center}

\onslide<2>
\begin{center}
\begin{tikzpicture}
\tikzset{component/.style={draw, rectangle, rounded corners=3pt, ultra thick, color=rathaxesred}}
\node[component] (cnorm)    at (-3, 0)      {\textcolor{white}{\LARGE{CNorm}}} ;
\node[component] (rtxparse) at (2, 2.1)     {\textcolor{white}{\Large{RTX\_Parse}}} ;
\node[hidden] (rtxtree)     at (2, 0.7)     {\Large{RTX\_Tree}} ;
\node[hidden] (rtxlink)     at (2, -0.7)    {\Large{RTX\_Link}} ;
\node[hidden] (rtxtpl)      at (2, -2.1)    {\Large{RTX\_Tpl}} ;
\node[hidden] (tools)       at (-5, -2.1)   {\Large{Tools}} ;
\draw[componentarrow] (cnorm) to (rtxparse) ;
\end{tikzpicture}
\end{center}

\onslide<3>
\begin{center}
\begin{tikzpicture}
\node[component] (cnorm)    at (-3, 0)      {\textcolor{white}{\LARGE{CNorm}}} ;
\node[component] (rtxparse) at (2, 2.1)     {\textcolor{white}{\Large{RTX\_Parse}}} ;
\node[component] (rtxtree)  at (2, 0.7)     {\textcolor{white}{\Large{RTX\_Tree}}} ;
\node[hidden] (rtxlink)     at (2, -0.7)    {\Large{RTX\_Link}} ;
\node[hidden] (rtxtpl)      at (2, -2.1)    {\Large{RTX\_Tpl}} ;
\node[hidden] (tools)       at (-5, -2.1)   {\Large{Tools}} ;
\draw[componentarrow] (cnorm) to (rtxparse) ;
\draw[componentarrow] (cnorm) to (rtxtree) ;
\end{tikzpicture}
\end{center}

\onslide<4>
\begin{center}
\begin{tikzpicture}
\node[component] (cnorm)    at (-3, 0)      {\textcolor{white}{\LARGE{CNorm}}} ;
\node[component] (rtxparse) at (2, 2.1)     {\textcolor{white}{\Large{RTX\_Parse}}} ;
\node[component] (rtxtree)  at (2, 0.7)     {\textcolor{white}{\Large{RTX\_Tree}}} ;
\node[component] (rtxlink)  at (2, -0.7)    {\textcolor{white}{\Large{RTX\_Link}}} ;
\node[hidden] (rtxtpl)   at (2, -2.1)       {\Large{RTX\_Tpl}} ;
\node[hidden] (tools)    at (-5, -2.1)      {\Large{Tools}} ;
\draw[componentarrow] (cnorm) to (rtxparse) ;
\draw[componentarrow] (cnorm) to (rtxtree) ;
\draw[componentarrow] (cnorm) to (rtxlink) ;
\end{tikzpicture}
\end{center}

\onslide<5>
\begin{center}
\begin{tikzpicture}
\node[component] (cnorm)    at (-3, 0)      {\textcolor{white}{\LARGE{CNorm}}} ;
\node[component] (rtxparse) at (2, 2.1)     {\textcolor{white}{\Large{RTX\_Parse}}} ;
\node[component] (rtxtree)  at (2, 0.7)     {\textcolor{white}{\Large{RTX\_Tree}}} ;
\node[component] (rtxlink)  at (2, -0.7)    {\textcolor{white}{\Large{RTX\_Link}}} ;
\node[component] (rtxtpl)   at (2, -2.1)    {\textcolor{white}{\Large{RTX\_Tpl}}} ;
\node[hidden] (tools)       at (-5, -2.1)   {\Large{Tools}} ;
\draw[componentarrow] (cnorm) to (rtxparse) ;
\draw[componentarrow] (cnorm) to (rtxtree) ;
\draw[componentarrow] (cnorm) to (rtxlink) ;
\draw[componentarrow] (cnorm) to (rtxtpl) ;
\end{tikzpicture}
\end{center}

\onslide<6>
\begin{center}
\begin{tikzpicture}
\node[component] (cnorm)    at (-3, 0)      {\textcolor{white}{\LARGE{CNorm}}} ;
\node[component] (rtxparse) at (2, 2.1)     {\textcolor{white}{\Large{RTX\_Parse}}} ;
\node[component] (rtxtree)  at (2, 0.7)     {\textcolor{white}{\Large{RTX\_Tree}}} ;
\node[component] (rtxlink)  at (2, -0.7)    {\textcolor{white}{\Large{RTX\_Link}}} ;
\node[component] (rtxtpl)   at (2, -2.1)    {\textcolor{white}{\Large{RTX\_Tpl}}} ;
\node[component] (tools)    at (-5, -2.1)   {\textcolor{white}{\Large{Tools}}} ;
\draw[componentarrow] (cnorm) to (rtxparse) ;
\draw[componentarrow] (cnorm) to (rtxtree) ;
\draw[componentarrow] (cnorm) to (rtxlink) ;
\draw[componentarrow] (cnorm) to (rtxtpl) ;
\end{tikzpicture}
\end{center}

\end{overprint}

\end{frame}

\begin{frame}{Merci}
\begin{center}
\LARGE{\itshape{Questions ?}}
\end{center}
\end{frame}

\end{document}
